% unidades/01-presentaciones.tex
\chapter{自己紹介 — Presentaciones}
\unidadheader{01}{自己紹介}{Presentaciones}{Conocer a alguien por primera vez}

\begin{tcolorbox}[vocabbox, title={📌 Vocabulario Nuevo}]
\begin{tabularx}{\linewidth}{llX}
\toprule
\textbf{Japonés} & \textbf{Español} & \textbf{Tipo} \\
\midrule
\vocabitem{名前}{なまえ}{nombre}{sust.}
\vocabitem{学生}{がくせい}{estudiante}{sust.}
\vocabitem{先生}{せんせい}{maestro / profesor}{sust.}
\vocabitem{会社員}{かいしゃいん}{empleado de empresa}{sust.}
\vocabitem{〜人}{じん}{persona de [país]}{sufijo}
\vocabitem{〜語}{ご}{idioma de [país]}{sufijo}
\vocabitem{〜歳}{さい}{años de edad}{sufijo}
\vocabitem{友だち}{ともだち}{amigo/a}{sust.}
\vocabitem{国}{くに}{país}{sust.}
\vocabitem{仕事}{しごと}{trabajo / ocupación}{sust.}
\bottomrule
\end{tabularx}
\end{tcolorbox}

\subsection*{🗣️ Frases Clave}

\frase{はじめまして。}{Mucho gusto. (al conocer a alguien)}
\frase{\ruby{私}{わたし}は〜と\ruby{申}{もう}します。}
  {Me llamo ~. (formal)}
\frase{どうぞよろしくお\ruby{願}{ねが}いします。}
  {Encantado/a de conocerle.}
\frase{お\ruby{名前}{なまえ}は\ruby{何}{なん}ですか？}
  {¿Cuál es su nombre?}
\frase{〜から\ruby{来}{き}ました。}
  {Vengo de ~ / Soy de ~.}

\subsection*{⚙️ Gramática}

\patron{Afirmación: ser / identificar}
  {〜は〜です}
  {\ruby{私}{わたし}はアルドです。
  \ej{\ruby{私}{わたし}は\ruby{学生}{がくせい}です。}{Soy estudiante.}
  \ej{\ruby{田中}{たなか}さんは\ruby{先生}{せんせい}です。}{El Sr. Tanaka es maestro.}}

\patron{Negación: no ser}
  {〜じゃありません \quad / \quad 〜ではありません}
  {じゃありません = informal \quad | \quad ではありません = formal
  \ej{\ruby{私}{わたし}は\ruby{会社員}{かいしゃいん}じゃありません。}{No soy empleado de empresa.}}

\patron{También: igual categoría}
  {〜も〜です}
  {\ej{\ruby{田中}{たなか}さんも\ruby{学生}{がくせい}です。}{El Sr. Tanaka también es estudiante.}}

\begin{tcolorbox}[ejerciciobox, title={📝 Ejercicios Rápidos}]
\begin{enumerate}
  \item ¿Cómo te presentas en japonés? (Di tu nombre con はじめまして)
  \item Traduce: \ruby{私}{わたし}はメキシコ\ruby{人}{じん}です。
  \item Niega la siguiente oración: \ruby{田中}{たなか}さんは\ruby{学生}{がくせい}です。
\end{enumerate}
\vspace{0.4em}
\textbf{Respuestas:}
1) はじめまして。[nombre]と申します。どうぞよろしく。\quad
2) Soy mexicano/a.\quad
3) \ruby{田中}{たなか}さんは\ruby{学生}{がくせい}じゃありません。
\end{tcolorbox}

\begin{tcolorbox}[kanjibox, title={🀄 Kanjis de la Unidad}]
\begin{tabularx}{\linewidth}{cllXX}
\toprule
\textbf{Kanji} & \textbf{On} & \textbf{Kun} & \textbf{Significado} & \textbf{Vocab} \\
\midrule
\kanjirow{名}{めい・みょう}{な}{nombre}{\ruby{名前}{なまえ} / \ruby{名刺}{めいし}}
\kanjirow{前}{ぜん}{まえ}{frente / antes}{\ruby{名前}{なまえ} / \ruby{前}{まえ}}
\kanjirow{学}{がく}{まな}{estudiar / aprender}{\ruby{学生}{がくせい} / \ruby{学校}{がっこう}}
\kanjirow{生}{せい・しょう}{い・う}{vida / nacer}{\ruby{先生}{せんせい} / \ruby{学生}{がくせい}}
\kanjirow{先}{せん}{さき}{antes / punta}{\ruby{先生}{せんせい} / \ruby{先}{さき}}
\bottomrule
\end{tabularx}
\end{tcolorbox}
